\section{Introduction}\label{sec:intro}

Let $\A$ be a finite alphabet and $\sh$ the left shift on $\A^{\Z}$. A reversible cellular automaton $F$ on $\A^{\Z}$ determines a two-dimensional subshift, its \emph{spacetime} $X_F\subset\A^{\Z^2}$, whose rows are the configurations $F^n(c)$, $n\in\Z$. The automorphism group of $X_F$ is the centraliser of $F$ in the group $\Aut(\A^{\Z})$ of invertible block maps (Lemma~\ref{lem:spacetime}): its elements are the invertible cellular automata that commute with $F$, that is, the invertible maps given by a translation-invariant local rule that commute with the dynamics. It always contains $\langle\sh,F\rangle$. We determine it for two nonlinear automata, with opposite outcomes, and obtain partial results for a third. The second of the two, $95$R, consists of two interleaved copies of the reversible automaton of Bobenko, Bordemann, Gunn and Pinkall \cite{BBGP1993}, known as Rule~54, and the group of the invertible cellular automata that commute with Rule~54 is determined as well (Theorem~\ref{thm:introB}(0)).

The automata are second-order elementary automata in the sense of Fredkin and Wolfram \cite{TM1990,Wolfram2002}. For an elementary rule $f\colon\{0,1\}^3\to\{0,1\}$ define, on $\A=\{0,1\}^2$ with cells written $(x,y)$,
\begin{equation}\label{eq:F}
F_f(c)_i=\bigl(y_i,\;x_i\oplus f(y_{i-1},y_i,y_{i+1})\bigr).
\end{equation}
Reading $y$ as the current row and $x$ as the previous one, this is the recursion $y_{t+1}=y_{t-1}\oplus f(y_t)$, reversible for every $f$. Following Takesue \cite{Takesue1987}, one writes $f$R for $F_f$. As models of statistical mechanics these automata were introduced by Takesue \cite{Takesue1987,Takesue1989}, who reduced them, by reflection and the cellwise complement, to $88$ classes \cite[Table~I]{Takesue1989}. We consider the rules $91$, $94$ and $95$, and their conjugates $218$, $122$ and $250$ under the cellwise complement. For these rules $f(a,b,c)=a\oplus c\oplus N(a,b,c)$ with a function $N$ that vanishes on most triples, so that the rows of a spacetime diagram satisfy a discrete wave equation with a source, and the fixed points of $\sh F_f$ and of $\sh^{-1}F_f$, the configurations that move at speed $1$ to the right and to the left, form two copies of a shift of finite type $\Phi_f$ of positive entropy (Section~\ref{sec:fixed}).

\subsection{Results}

\begin{thmx}\label{thm:introA}
$\Aut(X_{F_{94}})=\langle\sh,F_{94}\rangle\cong\Z^2$, and the same holds for rule $122$.
\end{thmx}

This is Theorem~\ref{cl:thm:dens}. The automaton $122$R is the example with which irreversible behaviour of a reversible rule is illustrated in \cite[p.~442]{Wolfram2002}. For the second theorem let $\mathsf G\subset\{0,1\}^{\Z}$ be the golden mean shift, the set of sequences with no two consecutive $1$s, and let $\mathcal N$ be the group of those automorphisms of $(\mathsf G,\sh)$ that preserve the number of $1$s of every point with finitely many $1$s. Let $\mathcal W$ be the set of all $u\colon\Z^2\to\{0,1\}$ with $u(j,k)\oplus u(j-1,k-1)=u(j,k-1)\vee u(j-1,k)$ for all $(j,k)$; it is the spacetime of the Rule~54 automaton in light-cone coordinates (Section~\ref{sec:lit}), and $\Aut(\mathcal W)$ is the group of the homeomorphisms of $\mathcal W$ that commute with the translations of $\Z^2$.

\begin{thmx}\label{thm:introB}
There is an isomorphism
\[
\Aut(X_{F_{95}})\;\cong\;\bigl(\mathcal N\times\mathcal N\times\Z\bigr)\big/\bigl\langle(\sh^{-1},\sh,-2)\bigr\rangle ,
\]
under which $\sh$ is the class of $(\mathrm{id},\mathrm{id},1)$ and $F_{95}$ the class of $(\mathrm{id},\sh,-1)$. Moreover:
\begin{enumerate}
\item[(0)] For the spacetime $\mathcal W$ of the Rule~54 automaton, $\Aut(\mathcal W)\cong\mathcal N\times\mathcal N$, the two light-cone translations corresponding to $(\sh,\mathrm{id})$ and $(\mathrm{id},\sh)$ (Theorem~\ref{cl:thm:95W}); that is, the group of the invertible cellular automata that commute with Rule~54 is $\mathcal N\times\mathcal N$.
\item[(1)] $\bigoplus_{m\ge1}S_m$ embeds in $\mathcal N$; so $\Aut(X_{F_{95}})$ contains every finite group and is not virtually $\Z^n$ for any $n$.
\item[(2)] $\langle\sh,F_{95}\rangle\cong\Z^2$ is central in $\Aut(X_{F_{95}})$, with quotient $(\mathcal N/\langle\sh\rangle)^2$; in particular it has infinite index.
\item[(3)] $F_{95}$ has no blocking word, and for all $(k,p)\in\Z^2$ with $p\ge1$ and all $M\ge0$ the window graph $W^{(k,p)}_M$ of $\sh^kF_{95}^p$ is strongly connected; equivalently, $F_{95}$ has no eigenfield of any speed.
\item[(4)] \textup{(computer-assisted)} Every automorphism of $X_{F_{95}}$ of radius at most $6$ lies in $\langle\sh,F_{95}\rangle$, and exactly $36$ automorphisms of radius $7$ do not.
\end{enumerate}
The same holds for rule $250$.
\end{thmx}

The isomorphism is Theorem~\ref{cl:thm:95Aut}; (0) is Theorem~\ref{cl:thm:95W}, in which $\mathcal W$ is defined by \eqref{cl:eq:Whalf}; (1) is Lemma~\ref{cl:lem:95box} and Corollary~\ref{cl:cor:95G}, (2) is Corollary~\ref{cl:cor:95quot}, (3) is Propositions~\ref{cl:prop:95noblock} and~\ref{cl:prop:95eigen}, and (4) is Proposition~\ref{cl:prop:95small}; blocking words, window graphs and eigenfields are defined in Section~\ref{sec:blocking}, and the statements for $122$ and $250$ follow by Lemma~\ref{lem:classes} ($\kappa$ is a cellwise conjugacy commuting with $\sh$, so blocking words and the graphs $W^{(k,p)}_M$ are carried over as well). The automaton $95$R, in the form $250$R, is the reversible automaton known as Rule~54, taken on both sublattices of spacetime (Section~\ref{sec:lit}); so Theorem~\ref{thm:introB}(0) gives the automorphism group of the spacetime of that automaton, and the displayed isomorphism that of two interleaved copies.

In the constructions of large automorphism groups for second-order automata in \cite{PaperC}, the automorphisms outside $\langle\sh,F\rangle$ of rules that are not affine come from a blocking word \cite{Kurka1997,BT2000} with an autonomous degree of freedom next to it (a blocking word cuts the line into parts that do not influence each other, and a local rule can then act on a degree of freedom that the word confines), and the affine rules $90$ and $165$, which have no blocking word, show that linearity can take its place. One may ask whether, for a rule that is not affine, the absence of blocking words forces $\Aut(X_{F_f})$ to be virtually $\Z^2$. By (1) and (3) the answer is negative (Corollary~\ref{cl:cor:95dich}): $f_{95}(a,b,c)=1\oplus ac$ is not affine. Whether the class of $95$ is the only exception among the elementary rules is Question~1.6 of \cite{PaperC}. Part (4) shows that an exhaustive search at small radius does not detect such a group: the block maps of radius at most $4$ that commute with $F_{95}$ and are invertible are the maps $\sh^aF^b$ (Lemma~\ref{cl:lem:95cens} and Proposition~\ref{cl:prop:95census}, computer-assisted).

Both theorems rest on the following statement, which we call Lemma~L. Let $\Phi^+$ and $\Phi^-$ be the sets of fixed points of $\sh F_f$ and of $\sh^{-1}F_f$.

\begin{thmx}[Lemma L]\label{thm:introC}
Let $f\in\{91,94,95\}$. A block map of any radius that commutes with $F_f$ and fixes every point of $\Phi^+\cup\Phi^-$ is the identity. Consequently two block maps that commute with $F_f$ and agree on $\Phi^+\cup\Phi^-$ are equal, and the restriction $\Aut(X_{F_f})\to\Aut(\Phi^+,\sh)\times\Aut(\Phi^-,\sh)$ is injective.
\end{thmx}

This is Theorem~\ref{cl:thm:L} with Corollaries~\ref{cl:cor:res} and~\ref{cl:cor:95}. For rule $91$ it is the main result: $\Aut(X_{F_{91}})$ is not determined. We show that an automorphism maps the points of $\Phi^\pm$ with a single $0$ to points that satisfy the constraints of Proposition~\ref{cl:prop:91KR}, and why the arguments used for $94$ and $95$ stop there (Remark~\ref{cl:rem:91}); Section~\ref{cl:sec:IR} contains two further statements on the scattering for this rule. For rule $155$, which is of the same form with a source $N$ that is not symmetric under reflection, the method of proof of Lemma~L fails: a front carrying the pair $11$ makes the source fire at every step (Proposition~\ref{cl:prop:155}), while to the right of a perturbation the medium is transparent (Proposition~\ref{cl:prop:155W} and Corollary~\ref{cl:cor:155W}). Whether Lemma~L holds for $155$ is open.

\subsection{The method}

\emph{The wave equation.} For $f\in\{91,94,95\}$ write $f(a,b,c)=a\oplus c\oplus N(a,b,c)$; then $N$ is the indicator of $000$ for $91$, of $010$ for $94$, and of $\{000,010\}$ for $95$. The $y$-bits $Y(t,i)$ of the rows $F^t(c)$ satisfy
\[
Y(t+1,i)\oplus Y(t-1,i)\oplus Y(t,i-1)\oplus Y(t,i+1)=N\bigl(Y(t,i-1),Y(t,i),Y(t,i+1)\bigr),
\]
and we say that $Y$ \emph{fires} where the right side is $1$. In the light-cone coordinates $q=t-i$, $p=t+i$ the points of $\Phi^+$ are the solutions $\pi(q)$ and those of $\Phi^-$ the solutions $\xi(p)$, with $\pi,\xi\in\Phi_f$, the sequences in which no triple is in the support of $N$. On a quadrant bounded by two light rays a solution is determined by its values on the two rays, as in a Goursat problem, and it is the sum of the two boundary values and the parity of the number of firings in a rectangle (the d'Alembert formula, Lemma~\ref{cl:lem:dalembert}).

\emph{Lemma L.} Let $c$ differ in finitely many cells from a spatially periodic configuration without firings. Along the light rays of each family, the first position at which the source can fire moves outwards by at least $3$ in every two steps (Lemmas~\ref{cl:lem:climb} and~\ref{cl:lem:climb95}), and the product of the two rates exceeds $1$; so $Y_c$ fires only finitely often in positive time (Lemma~\ref{cl:lem:fire}), and every window of a late row $F^t(c)$ is a window of a point of $\Phi^+\cup\Phi^-$ (Proposition~\ref{cl:prop:sep}). A block map $h$ that commutes with $F_f$ and fixes $\Phi^+\cup\Phi^-$ therefore fixes $F^t(c)$ for large $t$, hence $c$, and such configurations are dense.

\emph{Scattering.} Let $P\in\Phi^+$ be placed to the left of $X\in\Phi^-$, both with eventually $2$-periodic diagrams. The diagram of the resulting configuration fires finitely often and then consists of a point $\Gamma\in\Phi^-$ on the left and a point $\Omega\in\Phi^+$ on the right, which move apart (Lemma~\ref{cl:lem:scatter}); we write $\mathcal S(P,X)=(\Gamma,\Omega)$. Time reversal inverts $\mathcal S$ (Proposition~\ref{cl:prop:dual}), and every block map $h$ that commutes with $F_f$ intertwines it: $\mathcal S(hP,hX)=(h\Gamma,h\Omega)$ (Lemma~\ref{cl:lem:inter}).

\emph{Kink rigidity.} For $94$ the sequences of $\Phi_{94}$ are those without $010$, and the simplest non-periodic points of $\Phi^\pm$ are the \emph{bare kinks} $\cdots000111\cdots$ and $\cdots111000\cdots$. Two kinks that cross with the value $1$ between them and at equal parity make the source fire twice, and both are delayed by one step; otherwise they cross freely (Lemma~\ref{cl:lem:pair}). From these crossings and the intertwining property we show that every automorphism translates the bare kinks of $\Phi^+$ by one amount and those of $\Phi^-$ by another, of the same parity, so that after composition with an element of $\langle\sh,F\rangle$ it fixes all of them (Theorem~\ref{cl:thm:KR94}). The same holds for $95$ with the \emph{bare defects} $\cdots111\,0\,111\cdots$ in place of the kinks (Theorem~\ref{cl:thm:KR95}).

\emph{The window property and prepending, for $94$.} Let $g$ be an automorphism of radius $R$ that fixes the bare kinks. Its set of fixed points in $\Phi^+\cup\Phi^-$ is closed under $\mathcal S$ and $\mathcal S^{-1}$, and it contains every point of $\Phi^+$ all of whose words of length $2R+2$ are words of points of $\Phi^+$ that $g$ fixes, and likewise for $\Phi^-$: the \emph{window property}. The \emph{scattering closure} $C_L$ is the smallest set that contains the two constant configurations and the bare kinks, is invariant under the shift, the reflection and the time reversal, and has these properties for windows of length $L$ (Definition~\ref{cl:def:CL} and Lemma~\ref{cl:lem:closure}). A kink is prepended to a point $U$ of $C_L$ as follows (Lemma~\ref{cl:lem:prepend}): bare kinks are scattered off the point $[b;U]$ with the additional kink in such a way that the new kink is never delayed, while the first kink of the rest is delayed at every second scattering; after $2L$ scatterings the two are so far apart that the window property applies, and the closure under $\mathcal S^{-1}$ leads back to $[b;U]$. With a similar lemma for isolated zeros this shows that $C_L$ contains every point with finitely many kinks (Theorem~\ref{cl:thm:objects}), hence, by the window property, all of $\Phi^+\cup\Phi^-$, and Lemma~L gives $g=\mathrm{id}$.

\emph{The particle description, for $95$.} For $95$ the scattering closure of the bare defects is a proper subset of $\Phi^+\cup\Phi^-$, because $\mathcal S$ only translates the zeros of a point (Corollary~\ref{cl:cor:95trans}), and a normalised automorphism need not be the identity. Since $f_{95}$ does not read its middle argument, the two sublattices of a diagram evolve independently, each as a copy of a $\Z^2$-subshift $\mathcal W$ (Lemma~\ref{cl:lem:95split}), the spacetime of Rule~54. We use the known particle picture of Rule~54 \cite{BBGP1993} in the following form (Theorem~\ref{cl:thm:95sol}; see Section~\ref{sec:lit}): $\mathcal W$ is the image of $\mathsf G\times\mathsf G$ under a continuous map $\Psi$, at most two-to-one, with $\Psi(\sh^ax,\sh^by)=\tau^{\theta_{x,y}(a,b)}\Psi(x,y)$ for a cocycle $\theta$. The $1$s of $x$ and of $y$ are right- and left-moving particles, and the delays at their crossings are removed by this change of coordinates. An element $\varphi$ of $\mathcal N$ has a flux $\lambda_\varphi$ (Lemma~\ref{cl:lem:95flux}), and for $\varphi,\psi\in\mathcal N$ the map $\Psi(x,y)\mapsto\Psi(\varphi x,\psi y)$, corrected by the translation by $(\lambda_\psi(y),\lambda_\varphi(x))$, is a well-defined automorphism of $\mathcal W$; by Lemma~L these are all, so $\Aut(\mathcal W)\cong\mathcal N\times\mathcal N$ (Theorem~\ref{cl:thm:95W}). The group $\Aut(X_{F_{95}})$ is generated by the diagonal action of $\Aut(\mathcal W)$ on the two sublattices and by $\sh$, which exchanges them (Theorem~\ref{cl:thm:95Aut}); the proof uses Theorem~\ref{cl:thm:KR95} and Lemma~L.

\subsection{Relation to the literature}\label{sec:lit}

\emph{Rule 54.} On one sublattice the relation of $95$R is, after complementation, the update rule $s_x^{t+1}=s_x^{t-1}\oplus(s_{x-1}^t\vee s_{x+1}^t)$ on a checkerboard lattice: the system $\mathcal W$ is the spacetime of the reversible automaton introduced by Bobenko, Bordemann, Gunn and Pinkall \cite{BBGP1993}, known as Rule~54 or RCA54 and much studied as an exactly solvable model \cite{BKP2021}. In \cite{BBGP1993} the rule is written on a diagonally oriented lattice, as it is in the light-cone coordinates used here. That the spacetime pattern of a rule that does not read its middle argument separates into two independent parts, like the two colours of a checkerboard, is noted in \cite[p.~386]{Takesue1989}, where the classes of these rules are listed, among them that of $95$R, and again for $90$R, $95$R and $165$R in \cite[Section~4]{Takesue1995}; that $250$R decomposes into two independent copies of RCA54 is shown in \cite[Section~3.1]{IT2018}. (RCA54 is not the second-order automaton $54$R, whose automorphism group also contains every finite group, for a different reason \cite{PaperC}.)

The particle picture of Rule~54 is known, and we do not claim it: left- and right-moving particles at light speed whose only interaction is a delay by one step at each crossing. The conservation of the particles and their pairwise interaction are proved in \cite[Theorem~1, Corollary~1]{BBGP1993} for arbitrary configurations, by an inspection of cases; the delays are derived there for finitely many particles on the background $0$ \cite[Corollary~2]{BBGP1993}. Local rules that read off the particles of both families from an arbitrary configuration are given in \cite{GHKV2018,ADM2019,FGV2019,KlobasThesis}. The description of a state by the asymptotic spacings of its particles, with a hard-core constraint, is given in \cite{Gopalakrishnan2018}, where it is inferred from simulations, and is used in \cite{FGV2019}: in the asymptotic coordinates no two particles of one family are adjacent, while on the lattice two particles of one family are separated by at least two sites or have a particle of the other family between them \cite{FGV2019}, \cite[Section~6.1]{BKP2021}. None of these sources states or proves the global parametrisation: that every point of $\mathcal W$ is $u_{x,y}$ for some $(x,y)\in\mathsf G\times\mathsf G$, with the fibres of Theorem~\ref{cl:thm:95sol}(d) and the action of the translations of (b). The proofs of Theorems~\ref{cl:thm:95W} and~\ref{cl:thm:95Aut} need exactly this, and since no proof in this form is available in the sources, Theorem~\ref{cl:thm:95sol} is proved here. A rigorous description of another soliton automaton in coordinates in which the solitons move freely is the soliton decomposition of the box-ball system \cite{FNRW2021}.

Local conserved quantities of Rule~54 whose number grows with the size of the support like the Fibonacci numbers, identified with arrangements of particles that move in one direction, are reported from numerical searches in \cite[Section~2.2.3]{KlobasThesis} and \cite[Section~3.3]{BKP2021}, with the remark in \cite{KlobasThesis} that a systematic construction of these quantities is missing; because of them the model is called superintegrable in \cite{GP2022,GP2024,PP2026}. Operators that commute with the dynamics are known in the quantum setting: local Hamiltonians \cite{FGV2019,GP2024} and transfer matrices \cite{GP2024}, and, for quantum deformations of the model, charges that commute with the deformed dynamics \cite{PP2026}; Rule~54 is not among the Yang--Baxter integrable models with three-site interactions \cite{GP2021}. These are quantum operators and not cellular automata. Theorems~\ref{cl:thm:95sol} and~\ref{cl:thm:95W} are a rigorous counterpart of these observations within cellular automata: the arrangements of the particles of one family are the points of the golden mean shift $\mathsf G$, and every automorphism of $\mathsf G$ that conserves the number of particles, applied to one family, is an invertible cellular automaton that commutes with the dynamics. We have not found the group of the invertible cellular automata that commute with the dynamics (Theorems~\ref{cl:thm:95W} and~\ref{cl:thm:95Aut}), or the group $\mathcal N$, in the literature.

\emph{Conserved quantities.} The additive conserved quantities of the second-order elementary automata are determined for densities on two adjacent cells in \cite[Table~II]{Takesue1989} and on at most three adjacent cells in \cite[Tables~3--5]{HT1991}, and counted for densities on at most four adjacent cells in \cite[Table~1]{Takesue1997}; see also \cite{Takesue1995}. The two that we use for $91$ are recalled in Remark~\ref{cl:rem:91}; for one of them no verification is written out in the sources, and one is included (Lemma~\ref{cl:lem:91Q1}). That a cellular automaton on a full shift conserves an additive quantity if and only if the change of its density is the difference of a flux and its translate is known \cite[Theorems~2.1, 2.2, 2.6 and~2.7]{HT1991}, \cite{Pivato2002,MBG2004} (for number-conserving rules see also \cite{BF1998}); the framework on strongly irreducible shifts of finite type is in \cite[Section~4.1]{KT2015}, where the implication from conserved to locally conserved is stated for full shifts. Lemma~\ref{cl:lem:95flux} is this criterion on the golden mean shift, with the cocycle identities that we need.

\emph{Automorphism groups.} The automorphism group of a mixing shift of finite type such as $\mathsf G$ contains every finite group and free groups \cite{BLR1988}; the subgroup $\mathcal N$ of Theorem~\ref{thm:introB} still contains $\bigoplus_mS_m$, and we do not know its structure (Section~\ref{sec:questions}). The general results on automorphism groups of $\Z^2$-subshifts do not apply to the spacetime $X_F$: the positive-entropy theorem of Hochman \cite{Hochman2010}, because $X_F$ has zero $\Z^2$-entropy; the methods for one-dimensional subshifts of low complexity \cite{CK2015,DDMP2016}, because the rows of $X_F$ are all of $\A^{\Z}$; and the algebraic rigidity of Kitchens and Schmidt \cite{KS2000}, because the rules of this paper are not affine. The spacetime of a general endomorphism is studied in \cite{CFK2019}, and the problem of finding the block maps that commute with a given one goes back to \cite{CHR1979}.

\emph{Gliders.} The scheme of the proof of Theorem~\ref{thm:introA} has a precedent in the main lemma of Kopra on glider automorphisms \cite[Lemma~1]{Kopra2020}, \cite[Lemma~5.7]{Kopra2022}. There $X$ is a subshift with a group $\mathcal G$ of automorphisms that has left- and right-going fleets of gliders, defined as the finite points on which an element of $\mathcal G$ acts as a power of the shift, such that every finite point is mapped by an element of $\mathcal G$ to a point all of whose windows of a given length are windows of fleets, and such that for a sequence of distances $N_m\to\infty$ there are elements of $\mathcal G$ that delay a right-going and a left-going glider that meet at distance exactly $N_m$ and not those at larger distances (in \cite{Kopra2020} the fleets are indexed by a set, and the condition is imposed along the edges of a connected bipartite graph); the conclusion is that the automorphisms that commute with $\mathcal G$ and preserve the orbit of the background point are the powers of $\sh$. The three steps are those of Lemma~L, kink rigidity and the closure: the fleets correspond to $\Phi^\pm$, the diffusion to Proposition~\ref{cl:prop:sep}, and the delays to Lemma~\ref{cl:lem:pair}. In \cite{Kopra2020,Kopra2022} the group $\mathcal G$ is constructed for the purpose, on mixing shifts of finite type and on conjugates of transitive sofic shifts, and the result is a finitary form of the theorem of Ryan \cite{Ryan1972} that the centre of the automorphism group of a full shift consists of the powers of the shift: two automorphisms whose common centraliser is $\langle\sh\rangle$ (for the full $4$-shift a finite set with this property is given in \cite[Theorem~2.2]{Salo2019}). Here the automaton $F$ is given and not constructed, the sets $\Phi^\pm$ are shifts of finite type of positive entropy with several backgrounds, the delays occur only at equal parity, and the conclusion concerns the centraliser of the single map $F$.

\emph{Companion papers.} The group $\Aut(X_{F_f})$ is $\langle\sh,F_f\rangle$ for the non-affine rules that are permutive in an outer variable \cite{PaperB} and for the classes of $110$, $14$, $46$ and $79$, and it is $\langle\sh,F_f\rangle\times\Z/2$ for the classes of $24$ and $189$ \cite{PaperA}; in the cases with group $\langle\sh,F_f\rangle$ an automorphism is first normalised by an element of $\langle\sh,F\rangle$ and then shown to be the identity, as in the proof of Theorem~\ref{thm:introA}, but the structures used are different: for $110$ the fixed points of $\sh^{\pm1}F$ are two constants and one orbit of kinks, and here they form a shift of finite type of positive entropy. Large groups from walls and from linearity are constructed in \cite{PaperC}, and \cite{PaperD} records what is known for each of the $88$ classes. No result of the present paper depends on these papers.

\subsection{Computer-assisted statements}

The following statements rest on computations, and each is marked as such where it is stated: the enumeration of the particle-conserving block maps of $\mathsf G$ that depend on at most $8$ cells (Lemma~\ref{cl:lem:95enum}) and Proposition~\ref{cl:prop:95small}, that is, Theorem~\ref{thm:introB}(4); the census of the block maps of radius at most $4$ that commute with $F_{95}$ (Lemma~\ref{cl:lem:95cens}) and their identification (Proposition~\ref{cl:prop:95census}); the search for $175$ (Remark~\ref{cl:rem:175}); the computations reported in Remarks~\ref{cl:rem:91} and~\ref{cl:rem:IR} for $91$ and $94$; and those reported in Remark~\ref{cl:rem:155} and Appendix~\ref{app:155} for $155$. All of them are exact and exhaustive over a finite set. In Proposition~\ref{cl:prop:95small}(b), that the explicit block map $h_\ast$ of radius $7$ is an automorphism of $X_{F_{95}}$ outside $\langle\sh,F_{95}\rangle$ is verified on all windows from its tables, without the theory of Section~\ref{cl:sec:95}. Everything else is proved by hand, in particular Theorem~\ref{thm:introA}, Theorem~\ref{thm:introB} with (1)--(3), and Theorem~\ref{thm:introC}. Appendix~\ref{app:comp} describes the computations, and in addition computations that confirm results proved by hand; these are not used in any proof.

\medskip
\noindent\textbf{Organisation.} Section~\ref{sec:prelim} fixes notation, describes the fixed points of $\sh^{\pm1}F_f$, and defines blocking words, window graphs and eigenfields. Section~\ref{cl:sec:light} proves Lemma~L. Section~\ref{cl:sec:scatter} introduces the scattering map and proves the kink rigidity theorems for $94$ and $95$, and Section~\ref{cl:sec:94} proves Theorem~\ref{thm:introA}. Section~\ref{cl:sec:95} treats rule $95$ and proves Theorem~\ref{thm:introB}. Sections~\ref{cl:sec:91} and~\ref{cl:sec:155} contain the results on $91$ and $155$, and Section~\ref{sec:questions} lists questions. Appendix~\ref{app:comp} describes the computations, and Appendix~\ref{app:155} the front strips of $155$R.
